%=====================================================================
\chapter{Controlled Experiments in Software Engineering and HCI}\label{ch:experiments}
%=====================================================================
\locator{3}{Part~3 --- Research Designs for Computing}

\begin{outcomes}
By the end of this chapter you will be able to:
\begin{tight}
  \item \textbf{State} the three features that make a study an experiment.
  \item \textbf{Explain} what randomisation buys that no other technique can.
  \item \textbf{Choose} between between-subjects, within-subjects and factorial
    designs, and name what each costs.
  \item \textbf{Counterbalance} a within-subjects design and say which threat
    it removes.
  \item \textbf{Judge} the student-versus-professional question on evidence
    rather than reflex.
  \item \textbf{Specify} an experiment completely enough to be pre-registered.
\end{tight}
\end{outcomes}

\begin{prereq}
\chref{questions} for causal questions and variable roles, and
\chref{philosophy} for why this design belongs to the post-positivist
tradition. \chref{sampling} follows immediately: an experiment you cannot
power is not an experiment you can run.
\end{prereq}

\begin{keyterms}
experiment $\bullet$ manipulation $\bullet$ randomisation $\bullet$ control
$\bullet$ experimental unit $\bullet$ subject and object $\bullet$ treatment
$\bullet$ factor and level $\bullet$ between-subjects $\bullet$
within-subjects $\bullet$ crossover $\bullet$ carryover effect $\bullet$
counterbalancing $\bullet$ Latin square $\bullet$ blocking $\bullet$ factorial
design $\bullet$ interaction $\bullet$ pilot study
\end{keyterms}

\section{Three Features, and Only the Third Is Optional}

\begin{definitionbox}[Controlled experiment]
A study in which the researcher (1) \term{manipulates} at least one independent
variable, (2) \term{randomly assigns} experimental units to treatments, and
(3) \term{controls} other variables that could affect the outcome.

\medskip
Manipulation and randomisation are what license a causal conclusion. Control
improves precision, and where randomisation is absent but manipulation remains,
you have a \term{quasi-experiment} (\chref{quasiexp}) --- weaker, often the only
option, and requiring an explicit argument that assignment was not confounded.
\end{definitionbox}

\begin{conceptbox}[What randomisation actually buys]
This is the most misunderstood idea in empirical research, and it is worth
being precise.

\medskip
Randomisation does not make groups \emph{equal}. With 20 subjects per group
they will differ on experience, on the day they slept badly, on everything.
What randomisation does is make the assignment \textbf{independent of every
variable, including the ones you have not thought of and cannot measure}.

\medskip
That is the whole trick, and nothing else achieves it. Matching handles
confounders you have identified. Statistical adjustment handles confounders you
have measured. Only randomisation handles the confounder you have never heard
of --- which, in a field as young as ours, is most of them. It converts an
unknown bias into known, quantifiable noise, and the noise is what the
statistics in \chref{inference} are built to handle.
\end{conceptbox}

\subsection{Vocabulary that prevents confusion}

\begin{center}
\small
\begin{tabular}{@{}L{3.0cm}L{11.3cm}@{}}
\toprule
Experimental unit & The thing assigned to a treatment. Randomise it, count it,
                    and analyse at this level. Getting this wrong is the most
                    common analysis error in the field \\
Subject           & The person taking part \\
Object            & The material worked on --- the code, the document, the
                    dataset \\
Factor            & An independent variable being manipulated \\
Level / treatment & A value of a factor \\
Response          & The measured outcome \\
\bottomrule
\end{tabular}
\end{center}

\begin{alertbox}[If you randomise teams, your sample size is teams]
Suppose you assign four teams of ten developers each to two conditions and
analyse 40 individual results as if independent. You have reported a sample
size of 40 when the design supports 4. Members of a team share a codebase, a
manager, a culture and a coffee machine; their outcomes are correlated.

\medskip
This inflates significance dramatically, and it is exactly the error that makes
underpowered studies look conclusive. The remedies are to analyse at the unit
of randomisation, or to model the nesting explicitly with a mixed-effects model
(\chref{regression}). Deciding which \emph{before} collecting data is part of
portfolio piece P5.
\end{alertbox}

\section{Three Designs, and the Price of Each}

\armfig{
\begin{tikzpicture}[
  grp/.style={rounded corners=2pt, draw=primary!55, fill=primary!7,
              text width=1.85cm, align=center, font=\scriptsize,
              minimum height=0.72cm},
  hdr/.style={font=\scriptsize\bfseries, text=primary},
  note/.style={font=\scriptsize\itshape, text=neutral, text width=3.6cm, align=center},
  ar/.style={-{Stealth[length=1.8mm]}, draw=secondary}]

  % between
  \node[hdr] at (1.0,1.5) {Between-subjects};
  \node[grp] (bA) at (0,0.72) {S1--S20\\ Treatment A};
  \node[grp] (bB) at (0,-0.3) {S21--S40\\ Treatment B};
  \node[note] at (1.0,-1.55) {No carryover.\\ Needs twice the subjects.\\
        Between-subject variance is noise};

  % within
  \node[hdr] at (6.0,1.5) {Within-subjects};
  \node[grp] (wA) at (5.0,0.72) {S1--S20\\ A then B};
  \node[grp] (wB) at (5.0,-0.3) {S1--S20\\ B then A};
  \node[note] at (6.0,-1.55) {Each subject is own control.\\
        Half the subjects.\\ Carryover must be counterbalanced};

  % factorial
  \node[hdr] at (11.0,1.5) {$2\times2$ factorial};
  \node[grp] (f1) at (10.1,0.72) {A1\,B1};
  \node[grp] (f2) at (12.1,0.72) {A2\,B1};
  \node[grp] (f3) at (10.1,-0.3) {A1\,B2};
  \node[grp] (f4) at (12.1,-0.3) {A2\,B2};
  \node[note] at (11.1,-1.55) {Two factors at once.\\
        Reveals interaction.\\ Cells shrink fast};

  \draw[ar] (2.6,0.2) -- (3.9,0.2);
  \draw[ar] (7.6,0.2) -- (8.9,0.2);
\end{tikzpicture}}
{Three experimental designs. The choice is almost always forced by how many
subjects you can actually recruit, which is why \chref{sampling} should be read
before this decision is fixed rather than after.}{expdesigns}

\begin{center}
\small
\begin{longtable}{@{}L{2.6cm}L{5.4cm}L{6.0cm}@{}}
\toprule
\textbf{Design} & \textbf{Use when} & \textbf{Principal threat} \\
\midrule
\endfirsthead
\toprule
\textbf{Design} & \textbf{Use when} & \textbf{Principal threat} \\
\midrule
\endhead
\bottomrule
\endfoot
Between-subjects
  & Treatments cannot be un-learned; subjects are plentiful
  & Individual differences become error variance, so power is low. Randomise
    well, and consider blocking on experience \\
Within-subjects (crossover)
  & Subjects are scarce --- almost always, in software engineering
  & \emph{Carryover}: having done task A changes performance on B. Also
    fatigue and learning. Counterbalance, and use different but equivalent
    objects \\
Factorial
  & You suspect two factors interact
  & Cells shrink: a $2\times2$ with 40 subjects has 10 per cell. Interactions
    need substantially more power than main effects \\
Blocked
  & A known nuisance variable dominates, e.g. experience
  & None specific; blocking on a variable that turns out irrelevant costs
    little, so block when in doubt \\
\bottomrule
\end{longtable}
\end{center}

\subsection{Counterbalancing}

In a within-subjects design every subject sees every treatment, so order
becomes a confound. Counterbalancing distributes order systematically.

\begin{center}
\small
\begin{tabular}{@{}L{2.8cm}L{11.5cm}@{}}
\toprule
Two treatments  & Half the subjects do A then B, half B then A. Order becomes a
                  factor you can test rather than a bias you carry \\
Three or more   & Full counterbalancing needs $k!$ orders. A \term{Latin
                  square} uses $k$ orders in which each treatment appears once
                  in each position \\
Objects too     & Rotate the objects (the code, the task) against the
                  treatments as well, or the effect of treatment is confounded
                  with the difficulty of the material \\
\bottomrule
\end{tabular}
\end{center}

\begin{worked}[A $3\times3$ Latin square for treatments and objects]
Three review techniques (T1, T2, T3), three code artefacts (O1, O2, O3), and
subjects assigned to three groups. Each group sees every technique and every
artefact, and no group sees the same pairing.

\medskip
\begin{center}
\small
\begin{tabular}{@{}L{2.4cm}ccc@{}}
\toprule
 & \textbf{Session 1} & \textbf{Session 2} & \textbf{Session 3} \\
\midrule
Group 1 & T1 on O1 & T2 on O2 & T3 on O3 \\
Group 2 & T2 on O3 & T3 on O1 & T1 on O2 \\
Group 3 & T3 on O2 & T1 on O3 & T2 on O1 \\
\bottomrule
\end{tabular}
\end{center}

\medskip
Each technique appears once per session position, so learning and fatigue are
balanced across techniques. Each technique meets each artefact exactly once, so
artefact difficulty does not favour any technique.

\medskip
What it does \emph{not} handle: a genuine interaction between technique and
artefact --- if T1 is uniquely suited to O1 --- which a Latin square averages
away rather than detects. Detecting that needs a full factorial design and
correspondingly more subjects.
\end{worked}

\section{Human Subjects: the Students Question}

Every software engineering experiment using students is challenged on the same
ground: professionals would behave differently. The question deserves a better
answer than reflex in either direction, and there is evidence to give
one~\cite{falessi2018empirical, sjoberg2005survey}.

\begin{center}
\small
\begin{tabular}{@{}L{4.4cm}L{9.9cm}@{}}
\toprule
\textbf{Position} & \textbf{The considered version} \\
\midrule
Students are acceptable
  & For tasks where the relevant skill is one students possess; for testing a
    theory rather than estimating an industrial effect; for early formative
    work. Graduate students in particular are often comparable to junior
    professionals \\
Students are not acceptable
  & Where the task depends on experience the population lacks --- architectural
    judgement, legacy system familiarity, organisational negotiation. Here
    students are not a noisy proxy, they are a different population \\
\bottomrule
\end{tabular}
\end{center}

\begin{conceptbox}[How to answer it in your thesis]
Do not claim students are equivalent to professionals. Do the following
instead, and the objection loses most of its force:

\begin{tightnum}
  \item \textbf{Characterise your subjects} rather than labelling them.
    Years of programming, industrial months, familiarity with the language and
    with the task type. ``Students'' spans a first-year undergraduate and a
    part-time MS scholar with six years in industry.
  \item \textbf{State the target population} the inference is to, explicitly.
  \item \textbf{Argue the specific task.} Is the skill it demands one your
    subjects have?
  \item \textbf{Put it in threats to external validity} (\chref{validity}) with
    a direction: would professionals plausibly show a larger or smaller effect,
    and why?
\end{tightnum}
\end{conceptbox}

\section{Pilot First. Always}

A pilot with three or four subjects, analysed exactly as the real study will
be, is the cheapest insurance in empirical research.

\begin{center}
\small
\begin{tabular}{@{}L{4.4cm}L{9.9cm}@{}}
\toprule
\textbf{What a pilot catches} & \textbf{What it costs you to find later} \\
\midrule
Instructions that are ambiguous & Every subject interpreted the task
                                  differently; the data is uninterpretable \\
A task too long or too short    & Floor or ceiling effects; no variance to
                                  analyse \\
Instrumentation that fails      & Timing not logged; the primary outcome is
                                  missing \\
An outcome measure with no variance & The whole study, wasted \\
An analysis script that cannot run on the data shape & Discovered at the worst
                                  possible moment \\
\bottomrule
\end{tabular}
\end{center}

\begin{alertbox}[Write and run the analysis script on pilot data]
Run your full analysis on the pilot data --- or on simulated data with the same
shape --- before collecting anything real. If the script produces a number, the
design and the analysis are compatible. If it errors, or produces something you
cannot interpret, you have found a design flaw for the price of an afternoon.

\medskip
This also forces the analysis to be specified in advance, which is precisely
what pre-registration requires (\chref{prereg}). Pilot data must not be pooled
into the final analysis.
\end{alertbox}

\begin{templatebox}[C7 --- Experiment design sheet]
\begin{tabular}{@{}L{4.0cm}L{9.5cm}@{}}
Research question / hypothesis & \fillline{9.2cm} \\[7pt]
Factor(s) and levels           & \fillline{9.2cm} \\[7pt]
Experimental unit              & \fillline{9.2cm} \\[7pt]
Subjects: population, N, recruitment & \fillline{9.2cm} \\[7pt]
Objects and how they are assigned & \fillline{9.2cm} \\[7pt]
Design (between / within / factorial) & \fillline{9.2cm} \\[7pt]
Randomisation procedure        & \fillline{9.2cm} \\[7pt]
Counterbalancing / blocking    & \fillline{9.2cm} \\[7pt]
Response variable(s) and instrument & \fillline{9.2cm} \\[7pt]
Expected effect size, and its source & \fillline{9.2cm} \\[7pt]
Power and required N (\chref{sampling}) & \fillline{9.2cm} \\[7pt]
Planned analysis               & \fillline{9.2cm} \\[7pt]
Exclusion rules, decided now   & \fillline{9.2cm} \\[7pt]
Ethics approval / consent      & \fillline{9.2cm} \\[7pt]
\end{tabular}
\end{templatebox}

\begin{pitfall}[Five ways experiments in our field go wrong]
\begin{tight}
  \item \textbf{Analysis at the wrong level.} Randomising teams, analysing
    individuals.
  \item \textbf{No power analysis.} Twelve subjects, an effect of $d = 0.4$,
    and a null result reported as ``no difference''.
  \item \textbf{An outdated or badly tuned baseline.} Your treatment beats a
    comparison nobody would use.
  \item \textbf{Optional stopping.} Collecting until the result is significant
    is a way of guaranteeing significance (\chref{prereg}).
  \item \textbf{Exclusions decided after seeing the data.} Removing outliers
    that inconveniently weaken the effect. Write the rule in advance and apply
    it blind to condition.
\end{tight}
\end{pitfall}

\begin{traditions}{the controlled experiment}
\tradEMP{Its home. Manipulation plus randomisation is the strongest available
warrant for a causal claim, and the price is artificiality.}
\tradDSR{Experiments serve as summative artificial evaluation of an artefact
(\chref{designscience}); the artefact is the treatment.}
\tradQUAL{Complements rather than competes: an experiment establishes that an
effect occurred, and a qualitative strand explains why participants behaved as
they did. This is a standard explanatory sequential design (\chref{mixed}).}
\tradFORM{Experiments do not settle formal questions, but they do tell you
whether an asymptotic result matters at realistic input sizes --- which is
often the more useful question.}
\end{traditions}

\begin{lab}[Design one experiment completely]
Take a causal question from your alignment table (\chref{questions}) and fill in
template C7 completely --- every row, no blanks.

\medskip
Then submit it to the two tests that catch most flawed designs:
\begin{tightnum}
  \item \textbf{The unit test.} What exactly is randomised? Is your stated N the
    number of those things?
  \item \textbf{The null test.} Write the sentence you would publish if the
    result were exactly zero. If that sentence is not publishable, the study is
    not worth running as designed --- and you have learned this before spending
    six months on it.
\end{tightnum}
\end{lab}

\begin{checkpoint}
\begin{tightnum}
  \item Randomisation does not equalise groups. What does it do, and why can no
    amount of statistical adjustment substitute for it?
  \item You have 16 professional developers and three techniques to compare.
    Which design, and what specific threat does it introduce?
  \item In the Latin square worked example, what is averaged away rather than
    detected, and what design would detect it?
  \item Why must outlier exclusion rules be fixed before data collection, even
    though the rule itself may be entirely reasonable?
\end{tightnum}
\end{checkpoint}

\begin{chaptersummary}
\begin{tight}
  \item Manipulation plus randomisation licenses a causal claim; control adds
    precision.
  \item Randomisation makes assignment independent of \emph{unknown}
    confounders. Nothing else does this.
  \item Analyse at the unit of randomisation, or model the nesting. Randomising
    teams and counting individuals inflates significance.
  \item Between-subjects costs subjects; within-subjects costs carryover and
    must be counterbalanced; factorial costs cell size but reveals interaction.
  \item Latin squares balance order and material; they average away
    treatment-by-object interactions rather than detecting them.
  \item Characterise subjects rather than labelling them, state the target
    population, and put the inference limit in threats with a direction.
  \item Pilot, and run the real analysis script on pilot or simulated data
    before collecting anything.
\end{tight}
\end{chaptersummary}

\begin{reviewq}
  \item Controlled experiments are high in internal validity and low in
    external validity. For a specific question in your area, design a two-study
    programme in which each study's weakness is covered by the other, and say
    what it would cost.
  \item The students-versus-professionals debate has run for two decades
    without resolution. Is it resolvable in principle? If yes, what evidence
    would settle it; if no, what should the field do instead?
  \item Optional stopping is prohibited in frequentist analysis but is
    acceptable under some sequential and Bayesian designs. Explain why the
    difference exists, and whether it offers a practical route for
    resource-limited software engineering research.
  \item An experiment finds no effect with adequate power. Describe how you
    would write this up so that it is publishable, and identify the features of
    the reviewing system that make that difficult.
\end{reviewq}
